\chapter{Sprint 38 --- Portal Order (logic + UI a11y) + Shell/layout re-fix + module app icon (2026-07-23)}

\section{Bối cảnh}

Phiên interactive: user chốt ``làm HẾT Dev-actionable trong tab 5. Issue List, kể cả issue đã
\emph{Ready for Retest} và \emph{Retest Failed}''. Chiến lược git = \textbf{batch theo nhóm}
(ledger/qa\_sync vẫn per-issue). Tổng \textbf{28 issue} chia 3 batch, cộng 1 fix phát sinh cuối phiên
(icon app phân hệ backend). Dev tối đa set \texttt{Ready for Retest} (QA Standard §6, không tự
\texttt{Done}).

\begin{itemize}
  \item \textbf{Batch A} --- order logic (6 issue \texttt{WJ-ORD-001/002/003/004/017/020}), nhánh
    \texttt{dev/2026-07-22-orderA-logic}, fix \texttt{bd55ef9}.
  \item \textbf{Batch B} --- order UI a11y + khung giờ + stepper (15 issue
    \texttt{WJ-ORD-005/006/009/010/011/012/013/014/015/018/019/021/022} +
    \texttt{FUNC-MOB-ORDER-005/006}), nhánh \texttt{dev/2026-07-22-orderB-ui}, fix \texttt{cd1025d}.
  \item \textbf{Batch C} --- shell/layout/typography (7 issue \texttt{UI-01/02/03/04/06},
    \texttt{UI-MOB-HOME-002}, \texttt{UI-PC-SHELL-001}), nhánh \texttt{dev/2026-07-23-shellC}, fix
    \texttt{2e2e70d}. Phần lớn là \textbf{re-fix} các issue Sprint 34/37 đã thử nhưng
    \emph{Retest Failed}.
\end{itemize}

\section{Phần A --- Order logic (Batch A)}

\begin{itemize}
  \item \textbf{WJ-ORD-001 (Critical) --- qty âm/0/không nguyên.} Validate kiểu số nguyên
    $\geq$ min\_qty ở controller; nhập $-1$/$0$/phân số $\to$ reject, không gửi request rác.
  \item \textbf{WJ-ORD-002/003 --- step \& xoá nguyên tử.} Tăng/giảm $\pm$\texttt{min\_qty} nguyên tử;
    giảm tới min = xoá dòng, không để lại line qty=0 mồ côi.
  \item \textbf{WJ-ORD-004 --- BFCache.} Trang giỏ khi back/forward không phục hồi state cũ lệch với
    server (thêm chống bfcache stale).
  \item \textbf{WJ-ORD-017/020 --- tìm kiếm unaccent.} Gõ không dấu vẫn khớp SP có dấu (unaccent).
\end{itemize}

\section{Phần B --- Order UI / a11y / khung giờ (Batch B)}

\begin{itemize}
  \item \textbf{WJ-ORD-005 --- H1 mobile.} Đổi ``Sản phẩm'' $\to$ ``Đặt hàng'' (khớp bottom-nav +
    heading PC); bộ đếm ``N SP'' cùng hàng title, cập nhật theo số SP sau tìm/lọc.
  \item \textbf{WJ-ORD-006 --- banner khung giờ thực.} Warnbar + nút Gửi theo \texttt{order\_window\_open}
    server tính theo tz store: trong giờ = xanh + enable; ngoài giờ = đỏ + nêu giờ mở lại + disable
    (server vẫn chặn \texttt{ORDER\_TIME\_CLOSED}).
  \item \textbf{WJ-ORD-009 --- cột Quy cách có điều kiện.} Chỉ render khi $\geq 1$ SP có
    \texttt{description\_ecommerce}; hết cột toàn dấu ``---''.
  \item \textbf{WJ-ORD-010 --- aria-label giỏ.} Nút Thêm/stepper/ô số lượng có accessible name kèm tên SP;
    nhóm stepper \texttt{role=group}, ô số lượng \texttt{aria-live=polite}.
  \item \textbf{WJ-ORD-011 --- vùng chạm.} PC step 24$\to$30, del 28$\to$34; mobile add 46$\times$32
    $\to$44$\times$40, cart step/del $\to$36$\times$36. (Giới hạn: row compact BA chốt min-height 58
    nên tối đa $\sim$40, chưa đủ 44$\times$44 mọi nút --- cần BA nới row nếu muốn.)
  \item \textbf{WJ-ORD-012 --- token CTA đậm.} Thêm \texttt{-{}-wujia-cta \#0F7CA8} (contrast trắng
    4.71:1, đạt AA) cho nút chữ-trắng-trên-nền; brand \texttt{\#28A9DF} (2.69:1) GIỮ cho
    accent/icon/badge --- không đổi nhận diện.
  \item \textbf{WJ-ORD-013 --- safe-area.} Thanh tổng + trang giỏ mobile thêm
    \texttt{padding-bottom env(safe-area-inset-bottom)}; khối ghi chú cuộn trong flow, không phình thanh.
  \item \textbf{WJ-ORD-014 --- first viewport detail.} Ảnh chi tiết SP mobile \texttt{max-height 240px}
    căn giữa $\to$ giá/CTA lên first viewport.
  \item \textbf{WJ-ORD-015 --- nhãn khung qua đêm.} \texttt{from$\geq$to} hiển thị
    ``10:00 hôm nay -- 04:00 ngày mai (UTC+7)''; trong ngày ``10:00 -- 15:00 (UTC+7)''.
  \item \textbf{WJ-ORD-018 --- cuộn ngang panel giỏ PC.} \texttt{overflow-x:auto} bảng 5 cột trong cột
    phải hẹp $\to$ không đẩy body cuộn ngang (1440$\times$900).
  \item \textbf{WJ-ORD-019 --- focus ring.} \texttt{outline 2px + offset 2px} cho search/select/nút;
    (Enter submit đã hoạt động native --- form GET + nút submit).
  \item \textbf{WJ-ORD-021 --- stepper trên list mobile.} qty=0 hiện ``Thêm''; qty$>$0 morph
    ``$-$ N $+$''; route \texttt{/cart/step} nhận \texttt{product\_id}; giảm ở min = xoá + toast.
  \item \textbf{WJ-ORD-022 --- phân biệt nhãn đếm.} Số DÒNG = ``N mặt hàng''; tổng QUANTITY =
    ``Tổng số lượng / N'' --- hết lẫn 2 con số cùng nhãn.
  \item \textbf{FUNC-MOB-ORDER-005 --- ghi chú đơn mobile.} Thêm textarea \texttt{portal\_note} như PC,
    autosave hook \texttt{/cart/note}, không bị realtime swap xoá khi đang gõ.
  \item \textbf{FUNC-MOB-ORDER-006 --- nút Tìm mobile.} Nút submit 44$\times$44 (icon search) $\to$ submit
    chắc chắn, không phụ thuộc Enter bàn phím ảo. (Giới hạn: chọn submit tường minh thay vì live-filter
    debounce --- nhẹ request; tách issue riêng nếu BA muốn live.)
\end{itemize}

\section{Phần C --- Shell/layout/typography (Batch C, re-fix)}

\begin{itemize}
  \item \textbf{UI-01 (High) --- Top Bar float-offset.} Selector cũ \texttt{.header-navbar.floating-nav}
    (spec 0,2,0) THUA cascade Vuexy \texttt{body.navbar-floating .header-navbar.floating-nav} (0,3,0) lúc
    runtime. Fix: nâng lên \texttt{body.navbar-floating .header-navbar.floating-nav.wujia-navbar} (0,4,1)
    + \texttt{position:fixed} tường minh $\to$ neo x=300/y=0/w=1620/h=72.
  \item \textbf{UI-02 (Low) --- language pill.} PC pill 118$\times$40 ép \texttt{!important} +
    \texttt{line-height:1} + \texttt{box-sizing} (thắng padding-y Bootstrap đẩy h=46). Template ĐÃ đúng
    (\texttt{lang==vi\_VN} $\to$ cờ VN + ``Việt Nam''). Control ghi ``English'' là ĐÚNG khi account
    test \texttt{lang=en\_US} --- data, không phải bug markup.
  \item \textbf{UI-03 (High) --- account pill.} 202$\times$52 + \texttt{margin-right:24px} $\to$ x=1694
    ($1920-24-202$). Role subtitle (Manager/Owner/Staff) đã có trong template.
  \item \textbf{UI-04 (Low) --- mobile action circle.} 44$\to$38px, radius 22$\to$19, gap 12$\to$16px
    (token \texttt{\_variables.css}). Blank Shell chốt 38px (trước đóng ở 44px).
  \item \textbf{UI-06 --- typography nhòe.} Thêm \texttt{-webkit-font-smoothing:antialiased} +
    \texttt{-moz-osx-font-smoothing:grayscale} + \texttt{text-rendering:optimizeLegibility} trên
    \texttt{html body} $\to$ hết ``nhòe/răng cưa'' (retest cũ Sprint 34/35 chỉ đổi font-family, chưa bật
    smoothing). Mở rộng override Inter + cap heading weight $\leq$700 (Vuexy đặt 800).
  \item \textbf{UI-MOB-HOME-002 --- title Home.} Token section cấp 1 = 18px / \texttt{line-height 24px} /
    weight 700 (đủ 3 thuộc tính). Retest cũ fail vì \textbf{thiếu line-height} (chỉ set font-size).
  \item \textbf{UI-PC-SHELL-001 --- sidebar 300px.} Ép \texttt{.main-menu} width =
    \texttt{-{}-wujia-sidebar-width} (300px) + \texttt{.content margin-left} = 300px ở $\geq$1200px $\to$
    đóng dải nền hở 40px (x=260--300). Token 300px đã khai báo nhưng chưa áp vào width/margin.
\end{itemize}

\section{Phần D --- Icon app phân hệ backend (fleet/franchise)}

User hỏi: đã bỏ ảnh vào \texttt{static/description/icon.png} mà app phân hệ trong menu backend vẫn không
có icon. \textbf{Root cause:} icon app-drawer backend lấy từ \texttt{web\_icon} trên \textbf{menuitem gốc},
KHÔNG từ \texttt{static/description/icon.png} (file đó chỉ nuôi thẻ Settings$\to$Apps). Cả
\texttt{menu\_wujia\_fleet\_root} và \texttt{menu\_wujia\_root} có \texttt{icon.png} 128$\times$128 trên đĩa
nhưng thiếu \texttt{web\_icon} $\to$ icon trống. Fix (\texttt{59c8086}): thêm
\texttt{web\_icon="<module>,static/description/icon.png"} vào 2 menuitem gốc; có hiệu lực khi deploy
\texttt{-u} (Odoo đọc file, đổ \texttt{web\_icon\_data}).

\section{Gì đã làm (nghiệp vụ)}

Trang Đặt hàng của khách nhượng quyền được siết chặt cả logic lẫn giao diện: không còn gửi được số lượng
âm/0/lẻ, tăng-giảm-xoá sản phẩm chính xác, tìm kiếm không dấu vẫn ra kết quả; banner khung giờ hiện đúng
trạng thái mở/đóng theo giờ từng cửa hàng và khoá nút Gửi khi ngoài giờ. Nút bấm to hơn, có nhãn cho trình
đọc màn hình, nút Gửi đủ tương phản chữ trắng (đạt chuẩn tiếp cận AA) mà vẫn giữ nguyên màu thương hiệu
\texttt{\#28A9DF}. Phần khung (header/sidebar/typography) PC được đưa về đúng lưới thiết kế: thanh trên và
sidebar hết lệch, chữ hết nhòe, tiêu đề đồng cỡ. Cuối cùng, icon của 2 phân hệ Fleet \& Franchise trong menu
backend đã hiện đúng.

\section{Lý do/Trade-off}

Phần lớn Batch C là \textbf{re-fix} issue \emph{Retest Failed} của Sprint 34/37: nguyên nhân gốc không phải
``chưa làm'' mà là làm thiếu 1 mảnh (UI-06 thiếu font-smoothing, UI-MOB-HOME-002 thiếu line-height, UI-01
thua cascade Vuexy vì specificity thấp). Bài học: override Vuexy phải \textbf{đo specificity runtime}, không
tin \texttt{!important} là đủ --- selector \texttt{body.navbar-floating ...} của vendor là 0,3,0 nên custom
phải $\geq$0,4,x. Các issue UI-01/03/PC-SHELL-001 phụ thuộc computed-style runtime nên \textbf{không xác minh
được headless} --- ledger ghi rõ cần user đo lại trên server 1920$\times$1080. Token CTA tách riêng
(\texttt{\#0F7CA8}) thay vì đổi brand toàn cục để không phá nhận diện \texttt{\#28A9DF} BA-lock.

\section{Verify}

Build \texttt{-u wujia\_portal\_sale,wujia\_portal\_layout,wujia\_portal\_base --stop-after-init --no-http}:
94 module loaded, registry loaded, \textbf{0 ERROR/CRITICAL/Traceback}. Build lại
\texttt{-u wujia\_fleet,wujia\_franchise} sau fix icon: RC=0, Odoo validate \texttt{web\_icon} path OK.
Brace-balance 5 file CSS: cân. 28 issue $\to$ \texttt{Ready for Retest} qua \texttt{qa\_sync --apply}
(verify 28/28). Tất cả merge vào \texttt{main} (\texttt{1d80ef0}) + push \texttt{origin/main}. Các issue
phụ thuộc computed-style (UI-01/03/PC-SHELL-001) chờ user đo trực quan trên server sau deploy \texttt{-u all}.
